% ch1_c §1.4–1.5 (书页 15–22 / p029–p036)
%--B-TAIL--
f.c.c.格子的布里渊区就是这个倒格子的Wigner–Seitz原胞。这个图形我们已经在图 4(c)中见过——一个兼有正方形面与六边形面的多面体。正方形面的中心对应于正格子诸立方晶面法线的方向；六边形面的中心则是诸对角面法线的方向。

还有一点颇有意思：f.c.c.结构是一种“密堆积”（close-packed）结构。依次叠合一层层(111)面就可以把它构造出来，每一层都是一个六边形网络，就像由刚性球排成的那样。这也是金属常常采取这种结构的原因之一：金属的内聚并不依赖于强方向性的键，而在很大程度上是一种体积效应。

\section{Bloch 定理}

我们已经学会了如何表示具有晶格周期性的函数。但对于一种物理理论来说这些还不够；我们还需要考虑结构的各种激发，它们会破坏严格的平移对称性。这类激发有若干种，其中最重要的是格波（lattice waves），即原子绕其平衡位置的振动；以及电子态（electron states），即电子在静态晶格的场中的运动。在第10章我们还将考虑自旋波（spin waves），即定域在晶体各原子上的那些自旋的激发。

所有这些激发都以一个动力学方程——或一个Schrödinger方程，或一个自旋交换哈密顿量——为特征，它们\emph{在晶格平移下不变}。于是，若 $\mathcal{V}(\mathbf{r})$ 是位于 $\mathbf{r}$ 处的电子所感受到的势，则对一切 $\mathbf{l}$ 都有 $\mathcal{V}(\mathbf{r}+\mathbf{l})=\mathcal{V}(\mathbf{r})$。电子的波函数 $\psi(\mathbf{r})$ 必须满足Schrödinger方程
\begin{equation*}
\left(-\frac{\hbar^{2}}{2m}\nabla^{2}+\mathcal{V}(\mathbf{r})-\mathcal{E}\right)\psi=0,
\tag{1.37}
\end{equation*}
在把作用于 $\psi$ 的算符中的 $\mathbf{r}$ 换成 $\mathbf{r}+\mathbf{l}$ 之后，这个方程保持不变。

在格波和自旋波的情形，形式体系要更复杂一些，但原理相同。设
\begin{equation*}
(u_1,\,u_2,\,\dots,\,u_{\mathbf{n}},\,\dots)
\end{equation*}
是位于 $1,2,\dots,\mathbf{n},\dots$ 各格点上的原子的位移（或自旋位移，或自旋偏离算符，或随便什么东西）。运动方程（或哈密顿量）将把这些变量（或算符）包含在内，其方式是：如果我们作一个晶格平移 $\mathbf{l}$——也就是说，把公式里原先出现 $u_{\mathbf{n}}$ 的地方统统换成 $u_{\mathbf{n}+\mathbf{l}}$——我们将看不出任何差别。换言之，晶格的所有晶胞都等价而且不可分辨。这颇像所谓“宇宙学原理”（cosmological principle）：无论我们从哪一点观察，宇宙看起来本质上都是一样的。

为了把这个性质形式化，我们采用如下记号。令 $\mathcal{H}(\mathbf{0})$ 和 $|0\rangle$ 表示晶格平移之前的哈密顿算符和波函数；令 $\mathcal{H}(\mathbf{l})$ 和 $|\mathbf{l}\rangle$ 表示平移 $\mathbf{l}$ 之后的同样一些数学对象。于是，对电子而言 $|0\rangle\equiv\psi(\mathbf{r})$，$|\mathbf{l}\rangle\equiv\psi(\mathbf{r}+\mathbf{l})$。对格波而言，若 $|0\rangle$ 表示位于 $\mathbf{n}$ 处的原子正经历某一特定运动的态，则 $|\mathbf{l}\rangle$ 就表示位于 $(\mathbf{n}+\mathbf{l})$ 处的原子正作同一运动的态。

平移不变性的表述便是
\begin{equation*}
\mathcal{H}(\mathbf{l})=\mathcal{H}(\mathbf{0}).
\tag{1.38}
\end{equation*}
现在，本征态必须满足如下形式的方程\footnote{在格波的经典运动方程的情形，频率 $\omega$ 的平方扮演能量的角色，而“哈密顿量”则是在我们假定所处理的是对平衡的小偏离时导出的那组联立线性微分方程的矩阵。参看 §2.1。}
\begin{equation*}
\mathcal{H}(\mathbf{0})\,|0\rangle=\mathcal{E}\,|0\rangle.
\tag{1.39}
\end{equation*}
这个方程与
\begin{equation*}
\mathcal{H}(\mathbf{l})\,|\mathbf{l}\rangle=\mathcal{E}\,|\mathbf{l}\rangle,
\tag{1.40}
\end{equation*}
完全相同，后者只不过是把算符与态函数中的所有变量作一次形式上的重新标记而已。但由式 (1.38)，这意味着
\begin{equation*}
\mathcal{H}(\mathbf{0})\,|\mathbf{l}\rangle=\mathcal{E}\,|\mathbf{l}\rangle,
\tag{1.41}
\end{equation*}
也就是说，态函数 $|\mathbf{l}\rangle$ 也是 $|0\rangle$ 所满足的那个方程的解。由于 $|\mathbf{l}\rangle$ 未必与 $|0\rangle$ 恒同，看上去我们似乎只要把碰巧找到的任何一个解加以平移，就能生成运动方程的数不清的解。当然，所有这些解在能量上都是简并的。

这显然是荒谬的；新解必定都以某种方式与我们原来的解等价。有两种情形需要考虑：

(i) 设 $|0\rangle$ 确实是\emph{非简并}的（这种情况实际上绝不会发生，因为一切晶格除平移不变性之外还具有镜面对称性）。那么唯一的可能是：每个新态 $|\mathbf{l}\rangle$ 都是 $|0\rangle$ 的某个倍数，从而在物理上与它不可分辨。必定存在某个数 $\lambda_1$，使得
\begin{equation*}
|\mathbf{a}_1\rangle=\lambda_1|0\rangle
\tag{1.42}
\end{equation*}
这是沿 $\mathbf{a}_1$ 方向走一步的结果。归一化要求
\begin{equation*}
|\lambda_1|^{2}=1,
\tag{1.43}
\end{equation*}
于是可以取
\begin{equation*}
\lambda_1=e^{ik_1},
\tag{1.44}
\end{equation*}
其中 $k_1$ 为实数。类似地，对其余两个基本方向上的单位平移，可以取
\begin{equation*}
|\mathbf{a}_2\rangle=e^{ik_2}|0\rangle;\qquad|\mathbf{a}_3\rangle=e^{ik_3}|0\rangle,
\tag{1.45}
\end{equation*}
于是一般平移必须导致如下关系
\begin{align*}
|\mathbf{l}\rangle&\equiv|l_1\mathbf{a}_1+l_2\mathbf{a}_2+l_3\mathbf{a}_3\rangle\notag\\
&=e^{ik_1}|(l_1-1)\mathbf{a}_1+l_2\mathbf{a}_2+l_3\mathbf{a}_3\rangle\notag\\
&=e^{ik_1 l_1}|0+l_2\mathbf{a}_2+l_3\mathbf{a}_3\rangle\notag\\
&=e^{i(k_1 l_1+k_2 l_2+k_3 l_3)}|0\rangle
\tag{1.46}
\end{align*}
这里我们沿 $\mathbf{a}_1$ 方向走 $l_1$ 步、沿 $\mathbf{a}_2$ 方向走 $l_2$ 步，如此等等，每走一步都乘上相应的因子。

我们按如下形式定义一个矢量 $\mathbf{k}$：
\begin{equation*}
\mathbf{k}=k_1\mathbf{b}_1+k_2\mathbf{b}_2+k_3\mathbf{b}_3,
\tag{1.47}
\end{equation*}
其中 $\mathbf{b}_1,\mathbf{b}_2,\mathbf{b}_3$ 是 $\mathbf{a}_1,\mathbf{a}_2,\mathbf{a}_3$ 的倒三基矢（reciprocal triad），如式 (1.19) 中那样。关系式 (1.46) 便成为
\begin{equation*}
|\mathbf{l}\rangle=e^{i\mathbf{k}\cdot\mathbf{l}}|0\rangle,
\tag{1.48}
\end{equation*}
这正是我们想要证明的结果：

\emph{对于任何一个满足Schrödinger方程（或其经典或量子的对应形式）的波函数/态函数，都存在一个矢量 $\mathbf{k}$，使得平移一个格矢 $\mathbf{l}$ 等价于乘上相位因子 $\exp(i\mathbf{k}\cdot\mathbf{l})$。}

每个不同的态函数可以有不同的波矢（wave-vector）$\mathbf{k}$，但每一个都必须满足这个条件。晶格具有平移不变性这个强的初始前提，对元激发的形式强加了一条强的限制条件。

但我们还需要在更一般的情形下给出证明。

(ii) 设 $|0\rangle$ 是简并的。为简单起见，设它是二重简并的，于是肯定存在两个能量相同而又彼此不同的态 $|0\rangle_1$ 和 $|0\rangle_2$。那么，晶格平移 $\mathbf{a}_1$ 至多能产生这两个态本身的线性组合。于是
\begin{equation*}
\left.\begin{aligned}
|\mathbf{a}_1\rangle_1&=\mathsf{T}_1^{11}\,|0\rangle_1+\mathsf{T}_1^{12}\,|0\rangle_2,\\
|\mathbf{a}_1\rangle_2&=\mathsf{T}_1^{21}\,|0\rangle_1+\mathsf{T}_1^{22}\,|0\rangle_2,
\end{aligned}\right\}
\tag{1.49}
\end{equation*}
诸数 $\mathsf{T}_1^{11}$ 等之所以这样书写，是因为它们是矩阵 $\mathsf{T}_1$ 的元素；为了保持归一化，$\mathsf{T}_1$ 必须是幺正矩阵（unitary matrix）。为简洁起见，把这对方程写成
\begin{equation*}
\left(\,|\mathbf{a}_1\rangle\right)=\mathsf{T}_1\left(\,|0\rangle\right).
\tag{1.50}
\end{equation*}
但 $|0\rangle_1$ 和 $|0\rangle_2$ 既然是简并的，就不是唯一定义的。我们原本可以从任何另外两个态出发，只要它们是这两个态的线性组合。例如，我们也许可以从
\begin{equation*}
\left.\begin{aligned}
|0\rangle_1&=\mathsf{S}^{11}\,|0\rangle_1+\mathsf{S}^{12}\,|0\rangle_2,\\
|0\rangle_2&=\mathsf{S}^{21}\,|0\rangle_1+\mathsf{S}^{22}\,|0\rangle_2,
\end{aligned}\right\}
\tag{1.51}
\end{equation*}
出发，其中诸数 $\mathsf{S}^{11}$ 等是另一个（任意的）幺正矩阵 $\mathsf{S}$ 的元素。

现在我们这样选取 $\mathsf{S}$，使矩阵 $\mathsf{S}\mathsf{T}_1\mathsf{S}^{-1}$ 成为对角矩阵。按标准的代数理论，这总是做得到的。例如，设它约化为如下形式
\begin{equation*}
\mathsf{S}\mathsf{T}_1\mathsf{S}^{-1}=
\begin{pmatrix}
e^{ik_1} & 0\\
0 & e^{ik_1'}
\end{pmatrix}.
\tag{1.52}
\end{equation*}
很容易验证，态 $|0\rangle_1$ 和 $|0\rangle_2$ 现在在平移之下的变换方式，恰好就像它们是非简并的一样，即如式 (1.42–1.45) 中那样：
\begin{equation*}
|\mathbf{a}_1\rangle_1=e^{ik_1}|0\rangle_1,\qquad|\mathbf{a}_1\rangle_2=e^{ik_1'}|0\rangle_2.
\tag{1.53}
\end{equation*}
态 $|0\rangle_1$ 有一个波矢，其在 $\mathbf{a}_1$ 方向上的分量为 $k_1$；态 $|0\rangle_2$ 对同一平移方向则有一个分量为 $k_1'$（一般说来与之不同）的波矢。

但是其他方向的平移又如何呢？考虑 $\mathbf{a}_2$。对初始态 $|0\rangle_1$ 和 $|0\rangle_2$，将有另一个矩阵 $\mathsf{T}_2$ 确定这一平移的结果，即像式 (1.49) 和 (1.50) 中那样符号式地写成
\begin{equation*}
\left(\,|\mathbf{a}_2\rangle\right)=\mathsf{T}_2\left(\,|0\rangle\right).
\tag{1.54}
\end{equation*}
现在我们也可以把这个矩阵 $\mathsf{T}_2$ 化成对角形式，正如式 (1.52) 中那样，只要选好一组适当的起始函数。乍一看，这似乎与我们用来使 $\mathsf{T}_1$ 对角化的那组函数不相容。但请考虑如下两种可供选择的约化：
\begin{equation*}
\left(\,|\mathbf{a}_1+\mathbf{a}_2\rangle\right)=
\left\{\begin{aligned}
\mathsf{T}_2(\,|\mathbf{a}_1\rangle)&=\mathsf{T}_2\mathsf{T}_1(\,|0\rangle),\\
\mathsf{T}_1(\,|\mathbf{a}_2\rangle)&=\mathsf{T}_1\mathsf{T}_2(\,|0\rangle).
\end{aligned}\right\}
\tag{1.55}
\end{equation*}
矩阵 $\mathsf{T}_1$ 和 $\mathsf{T}_2$ 由此彼此对易。矩阵代数中有一条定理告诉我们：这时存在一个幺正矩阵 $\mathsf{S}$，使得 $\mathsf{T}_1$ 和 $\mathsf{T}_2$ \emph{二者}都能\emph{同时}约化为对角形式。于是，态 $|0\rangle_1$ 和 $|0\rangle_2$ 不会被平移 $\mathbf{a}_2$ 混合，而只是被乘上相位因子
\begin{equation*}
|\mathbf{a}_2\rangle_1=e^{ik_2}|0\rangle_1,\qquad|\mathbf{a}_2\rangle_2=e^{ik_2'}|0\rangle_2.
\tag{1.56}
\end{equation*}
对平移 $\mathbf{a}_3$ 也有类似的结果，于是我们基本上回到了与式 (1.45) 实质相同的情形：对 $|0\rangle_1$ 和 $|0\rangle_2$ 中的每一个都存在一个波矢，使得
\begin{equation*}
|\mathbf{l}\rangle=e^{i\mathbf{k}\cdot\mathbf{l}}|0\rangle.
\tag{1.57}
\end{equation*}
把这一论证推广到两个以上能量简并的态，显然是轻而易举的。

这样，定理便在一般情形下得证，其含义是：波函数的任何一个简并解都可以表示成一些同能量解的线性组合，其中每一个都必须满足这类条件，尽管并非全都取同一个 $\mathbf{k}$ 值。

这就是\emph{Bloch定理}。其实，还存在一个一般的群论证明，它是下述定理的一条推论：“在复数域中，Abel群的任何表示都可以约化为若干一维表示之和”。晶体的平移群是Abel群——也就是说，群的所有元素彼此对易。这来自一个明显的事实：比如平移 $(\mathbf{a}_1+\mathbf{a}_2)$ 与 $(\mathbf{a}_2+\mathbf{a}_1)$ 是同一回事——我们可以沿许多不同的路径到达一个格点，但结果完全一样。矩阵 $\mathsf{T}_1$、$\mathsf{T}_2$ 等等正是这些平移的表示，因而可以约化为对角形式。

\section{约化到一个布里渊区}

Bloch定理的普遍性如此之强，以致在这个阶段我们很难把握它。对电子波而言，它意味着我们可以用波矢 $\mathbf{k}$ 来标记每一个波函数，并写成
\begin{equation*}
\psi_{\mathbf{k}}(\mathbf{r}+\mathbf{l})=e^{i\mathbf{k}\cdot\mathbf{l}}\psi_{\mathbf{k}}(\mathbf{r}).
\tag{1.58}
\end{equation*}
我们注意到，普通的自由电子波满足这个条件：
\begin{equation*}
\psi_{\mathbf{k}}(\mathbf{r})=e^{i\mathbf{k}\cdot\mathbf{r}}.
\tag{1.59}
\end{equation*}
这是意料之中的，因为我们处理的仍然是在周期势中Schrödinger方程的解；只不过这个势碰巧处处为零。空晶格检验（empty-lattice test）在固体理论中往往非常有用。

有时为了方便，人们试图让具有给定 $\mathbf{k}$ 值的电子波函数尽可能看起来像自由电子波。我们令
\begin{equation*}
\psi_{\mathbf{k}}(\mathbf{r})=e^{i\mathbf{k}\cdot\mathbf{r}}u_{\mathbf{k}}(\mathbf{r})
\tag{1.60}
\end{equation*}
并指望 $u_{\mathbf{k}}$ 将近乎恒定。事实上，由式 (1.58) 立刻可以证明，函数 $u_{\mathbf{k}}(\mathbf{r})$ 必须是周期的，即
\begin{equation*}
u_{\mathbf{k}}(\mathbf{r}+\mathbf{l})=u_{\mathbf{k}}(\mathbf{r}).
\tag{1.61}
\end{equation*}
Bloch定理有时就表述成这种形式。

定理中出现的因子 $\exp(i\mathbf{k}\cdot\mathbf{l})$，与我们研究周期函数时遇到的因子 $\exp(i\mathbf{g}\cdot\mathbf{l})$ 相类似。波矢 $\mathbf{k}$ 显然与倒格矢 $\mathbf{g}$ 具有相同的量纲；它属于倒格子空间。如果某个态碰巧具有波矢 $\mathbf{g}$，那它就应该是一个周期函数：
\begin{align*}
\psi_{\mathbf{g}}(\mathbf{r}+\mathbf{l})&=e^{i\mathbf{g}\cdot\mathbf{l}}\psi_{\mathbf{g}}(\mathbf{r})\notag\\
&=\psi_{\mathbf{g}}(\mathbf{r}),
\tag{1.62}
\end{align*}
因为
\begin{equation*}
e^{i\mathbf{g}\cdot\mathbf{l}}=1
\tag{1.63}
\end{equation*}
对一切 $\mathbf{l}$ 成立。

再设态 $\psi_{\mathbf{k}}$ 具有满足下式的波矢 $\mathbf{k}$：
\begin{equation*}
\mathbf{k}=\mathbf{g}+\mathbf{k}',
\tag{1.64}
\end{equation*}
其中 $\mathbf{g}$ 是某个倒格矢，而 $\mathbf{k}'$ 是另一个矢量。于是由式 (1.58) 和 (1.63)，
\begin{align*}
\psi_{\mathbf{k}}(\mathbf{r}+\mathbf{l})&=e^{i(\mathbf{g}+\mathbf{k}')\cdot\mathbf{l}}\psi_{\mathbf{k}}(\mathbf{r})\notag\\
&=e^{i\mathbf{g}\cdot\mathbf{l}}e^{i\mathbf{k}'\cdot\mathbf{l}}\psi_{\mathbf{k}}(\mathbf{r})=e^{i\mathbf{k}'\cdot\mathbf{l}}\psi_{\mathbf{k}}(\mathbf{r}).
\tag{1.65}
\end{align*}
这就是说，态 $\psi_{\mathbf{k}}$ 满足Bloch定理，仿佛它具有波矢 $\mathbf{k}'$ 似的。原来的标记 $\mathbf{k}$ 并不唯一；每一个态都拥有一大群可能的波矢，它们彼此之间相差倒格子的矢量。这当然不与定理矛盾，因为定理只不过断言每个态必须至少有\emph{一个}波矢。

这样我们就面临一个问题：如何唯一地定义一个给定态的波矢？可以拿式 (1.60) 作指导，试着选取 $\mathbf{k}$ 使 $u_{\mathbf{k}}(\mathbf{r})$ 尽可能接近常数——但这是任意的，而且正如我们以后会看到的，甚至会引人误入歧途。正确的做法如下：

考虑在一维情形（参看式 (1.8)–(1.12)）会发生些什么：倒格子的类似物是倒格子长度（reciprocal lattice lengths）的集合
\begin{equation*}
g_n=n\frac{2\pi}{a}.
\tag{1.66}
\end{equation*}
一个态可以被赋予集合
\begin{equation*}
k=n\frac{2\pi}{a}+k',
\tag{1.67}
\end{equation*}
中的任何一个波数；也就是说，$k$ 只是在模 $(2\pi/a)$ 的意义下确定的。图 12 中所有的 $\mathbf{k}$ 点都彼此等价。

%FIG{12}: {诸点 $k$ 在一维倒格子中全都约化为 $k'$。}

很自然地取 $k'$ 作为所有这些点的代表，且使 $|k'|$ 尽可能小。换句话说，我们总是选取落在范围
\begin{equation*}
-\frac{\pi}{a}<k\leqslant\frac{\pi}{a}.
\tag{1.68}
\end{equation*}
之内的值来作为波数。

显然，这个范围对我们的一维系统来说就等价于一个布里渊区；它是倒格子的一个晶胞。在三维情形我们如法炮制——我们选取波矢，使它落在倒格子空间的第一布里渊区内。这件事总是可能的，由初等几何即可明白。我们在式 (1.64) 中这样选取 $\mathbf{g}$ 的值，使 $|\mathbf{k}'|$ 尽可能小——也就是说，使它离倒格子的原点尽可能地近。这意味着 $|\mathbf{k}'|$ 所在之处离原点要比离倒格子的任何其他格点都近——这就等于说它落在该格子的Wigner–Seitz原胞内，也就是落在布里渊区内。显然，我们可以把倒格子空间中的任何一点 $\mathbf{k}$ \emph{约化}为布里渊区中的一点，于是任何态都可以在简约区图式（reduced zone scheme）中获得一个标记。

于是任何态都可以用它的简约波矢（reduced wave-vector）来表征。但可能有许多态具有同一个简约波矢而能量各不相同。因此，采用简约区图式使我们无法给每个态指定一个互不相同的 $\mathbf{k}$ 值。

%FIG{13}: {诸波矢 $\mathbf{k}$ 全都约化为位于布里渊区中的 $\mathbf{k}'$。}

看一看自由电子态在空晶格中会发生些什么，是很有意思的。设
\begin{align*}
\psi&=e^{i\mathbf{k}\cdot\mathbf{r}}\notag\\
&=e^{i(\mathbf{k}-\mathbf{g})\cdot\mathbf{r}}e^{i\mathbf{g}\cdot\mathbf{r}}\notag\\
&=e^{i\mathbf{k}'\cdot\mathbf{r}}\left(e^{i\mathbf{g}\cdot\mathbf{r}}\right),
\tag{1.69}
\end{align*}
其中 $\mathbf{k}'$ 是实际波矢 $\mathbf{k}$ 的约化值。这具有式 (1.60) 的形式，因为 $\exp(i\mathbf{g}\cdot\mathbf{r})$ 是晶格中的周期函数；但它是平面波的一种极其生硬人为的表示。如果在一维情形下能量由
\begin{equation*}
\mathcal{E}(\mathbf{k})=\frac{\hbar^{2}k^{2}}{2m},
\tag{1.70}
\end{equation*}
给出，那么在简约区中它将表现为 $\mathbf{k}'$ 的多值函数。在图 14 中，抛物线的 $AB$ 段经平移一个倒格矢被移到 $A'B'$ 处，如此等等。显然，在这种情形下，其中每个态都用它的“真实”波矢 $\mathbf{k}$ 来表示的扩展区图式（extended zone scheme）要更自然。
